\documentclass[sigconf,nonacm]{acmart}
\citestyle{authoryear}
\setcopyright{none}
\settopmatter{printacmref=false}
\renewcommand\footnotetextcopyrightpermission[1]{}

% Replace these placeholders with your name and affiliation.
\author{Andrew Chang}
\affiliation{%
  \institution{Carnegie Mellon University}
  \city{Pittsburgh}
  \country{USA}}

\begin{document}
\title{ITERATIVE SPECIFICATION REFINEMENT FOR ARCHITECTURAL ADHERENCE IN LLM-GENERATED CODE}
\maketitle

\section*{Questions for Reviewers}

\noindent\textbf{Section 1: Introduction}
\begin{itemize}
  \item Is it clear that reusable specification guidance is the contribution, and iterative refinement is how we develop it?
  \item Do the research questions match that contribution and the experiment's scope?
  \item Does the training-exposure paragraph explain what the comparison can establish?
\end{itemize}

\noindent\textbf{Section 4: Methods}
\begin{itemize}
  \item Are the reasons for choosing FEA-Bench, the evaluation criteria, and specification categories convincing?
  \item Is the generation--evaluation--revision loop clear enough to reproduce?
  \item Do the revision histories and simplification results provide appropriate evidence for the claims?
  \item Are any experimental decisions still stated without enough justification?
\end{itemize}

\section{Introduction}

In industrial software engineering, code added to production codebases must fit the structure of the existing codebase. While Large Language Models (LLMs) can generate functionally correct code, their implementations often don't align with the developer's intent. For example, LLM-generated code might pass unit tests but create a new class instead of using an existing one or implement a feature in the wrong module within the codebase. Developers then have to reprompt the model or refactor its output to fit the repository or their intended implementation. This creates additional work on maintaining LLM output quality and leaves developers feeling like it would have been easier to write code themselves.

Existing classes and interfaces may already handle validation, error handling, or interactions with other modules. Implementing the same behavior through a separate mechanism can duplicate these responsibilities or introduce dependencies that the repository was designed to avoid.

To implement a feature the way a developer envisions, an LLM must understand both the developer's intended behavior and the repository's implementation patterns. Developers can use software specifications to communicate these requirements, but it is unclear which specification content and formulations help an LLM follow the intended design. Long specifications may leave important implementation details unspecified while obscuring the intent of the feature. Shorter specifications risk excluding important details. For example, twenty broad user stories might describe the expected behavior without identifying an interface the feature must implement, whereas five refined, targeted requirements describing that interface, its registration mechanism, allowed dependencies, and input/output contract may provide more useful instructions.

Existing research examines what instructions developers provide to coding assistants and how clarifying requirements improves code generation. \citet{jiang2025} describe repository instructions, while ClarifyGPT demonstrates improvements from requirements clarification~\citep{mu2024}. These approaches inform this study, which examines how specifications can communicate repository-specific implementation expectations and how generated code can reveal what those specifications fail to communicate.

This research identifies what specification content  helps LLMs follow existing repository patterns when implementing new features. The paper explores an iterative specification refinement approach using reference code and evaluation feedback to improve LLM architectural adherence when adding features to existing codebases. The intended contribution is generalizable guidance on specification content and organization that improves architectural adherence. Reference-informed iterative refinement is the method used to develop that guidance; its practical value is evaluated on new tasks without reference-code input.

The scope is limited to feature additions within existing codebases. The goal is not to reproduce the reference implementation line by line, but to communicate the behavior, public contracts, and implementation constraints that the feature must satisfy. The agent retains freedom to make implementation decisions within these requirements.

The experiment uses historical feature requests and their merged pull requests from FEA-Bench~\citep{fea2025}. These tasks provide both the original codebase and a completed implementation that can serve as a reference. A code-generating agent receives a specification and the original codebase, without being supplied with the reference implementation during generation. Its output is evaluated against fixed behavioral and structural criteria using the reference code. A separate LLM acting as a specification editor uses the current specification, generated code, reference code, and evaluation feedback to revise the instructions. Each revision is tested through fresh code-generation runs until it consistently passes evaluation or reaches the revision limit.

It's possible that the features requests from FEA-Bench were used to train the LLM model used for this paper, GPT 6 Luna. We therefore evaluate baseline and refined specifications on the same tasks using the same model and original repository states. The comparison measures performance under different specification inputs, rather than performance on tasks guaranteed to be absent from training data.

The revision history provides more than a record of whether a specification passes. It connects implementation failures with changes to specification content and wording, showing what information was introduced or clarified before later implementations succeeded. These observations identify changes associated with improvement, rather than independently proving the effect of every instruction.

The research addresses three questions:

\begin{enumerate}
  \item Which specification content and wording are associated with correcting implementation failures during refinement?
  \item Can iterative specification refinement using reference code and evaluation feedback produce implementations that consistently pass the required behavioral and structural criteria?
  \item Does the guidance derived from these findings improve specification generation for new features and repositories?
\end{enumerate}

Each specification is treated as an executable hypothesis; its instructions are tested by generating code and evaluating the result. Running 10 of the same agents in parallel helps distinguish success from a single ``lucky'' attempt. The revision history records which instructions changed in response to failures and whether subsequent implementations passed evaluation. These observations inform guidance for generating specifications, which is then tested on unseen feature requests and repositories without access to their completed implementations. The practical goal is to help developers write specifications for new features, potentially through a reusable skill based on the research findings.

\section{Background}

The scope of this research is limited to brownfield development---features within an existing codebase. Because the codebase already has established design decisions, additional functionality must integrate with existing modules and follow the repository's conventions.

Problems can arise when an LLM implements the requested behavior without correctly identifying the interfaces and dependencies it must use. Interfaces and dependencies can span multiple files, and these conventions may only be demonstrated through the code rather than in documentation. \citet{tambon2025} shows the effect of this by identifying defects in LLM-generated code such as references to nonexistent objects, incorrect API usage, and misinterpreted requirements. These defects motivate the use of targeted specifications that clarify not only expected behavior but also implementation ambiguity. 

Following the repository's structure involves more than implementing the requested behavior. In this study, we defined structural coherency as whether generated code follows the repository's intended organization, including file placement, API contracts, dependencies, etc. An API contract specifies the names, parameters, types, return values, and behavior that other code relies on. An extension point is an existing mechanism for adding functionality, such as an interface or registration system. Creating a separate mechanism could implement the requested behavior but fail to follow the repository's intended design and conventions. Architectural adherence also includes following local conventions for validation, error handling, naming, logging, documentation, and tests.

Software specifications communicate these behavioral and structural requirements to the code-generating agent. Requirements engineering (RE) provides methods for identifying, documenting, and validating those requirements. The requirements engineering approach described by van Lamsweerde~\citep{vanlamsweerde2009} explains how high-level goals are refined into requirements, domain models, and software specifications. This refinement provides an example of prior work done on refining general feature descriptions into more specific instructions about data models, API contracts, and implementation constraints that can help an agent implement code with human-level decision making.

These requirements engineering practices also inform specification-driven development, an approach in which explicit software specifications guide planning, implementation, and verification. Specifications allow developers to describe the behavior the implementation must satisfy, along with requirements for how it integrates with existing code. The specification gives the LLM agent instructions for implementing the feature and provides expectations against which its output can be evaluated. By making these constraints explicit, specification-driven development can outline existing implementation patterns to limit the agency of an LLM. There is the added benefit that it forces both developers and LLMs alike to have a clear understanding of what constitutes a successful implementation.

The specification categories and structural evaluation criteria also draw on requirements engineering and architecture description standards. ISO/IEC/IEEE 29148:2018~\citep{iso29148} defines requirements as statements of needs and specifies what software requirements documents should contain, including expected behavior, performance requirements, external interfaces, and design constraints. ISO/IEC/IEEE 42010:2022~\citep{iso42010} explains how architectural views address stakeholder concerns. A view describes architecture from a perspective that helps examine the issue. For example, a view of component dependencies shows how changes might affect other components. This study builds on these concepts by applying them to the prompts for generating specifications and evaluating the defined set of criteria.

Evaluating whether these instructions produce the intended implementation requires a reference for both behavior and structure. In this study, the reference implementation is the completed code from a merged pull request, which provides the target behavior, public contracts, structure, etc. Evaluation checks whether generated code meets these expectations rather than whether it reproduces the reference patch exactly. This allows LLM agents to take agency in how it implements a feature bounded by a set of evaluation criteria that is important in brownfield development.

Different expectations require different forms of evaluation. Tests check specified behavior, while structural checks examine properties such as file placement, API signatures, and dependency directions. Whether responsibilities belong in the intended component or follow local conventions also requires subjectively interpreting the surrounding code. An LLM-as-a-judge assesses these criteria using the generated and reference code. Its judgments are checked against manual inspection because an incorrect judgment can misdirect the next specification revision.

\section{Related Work}

Architectural nonconformance requires manual correction even when code is functionality correct. This can make LLM agent usage more difficult than implementing features by hand. \citet{li2023} analyze 606 architecture-related review comments from four OpenStack and Qt projects. Refactoring and removing code account for 90\% of the measures used to address the identified violations. Although their study concerns human development rather than LLM-generated code, it shows that overall correctness isn't just determined by functional correctness but also architectural conformance.

Developers already communicate repository expectations through instructions for coding assistants. \citet{jiang2025} analyze cursor rules from 401 repositories, identifying conventions, guidelines, project information, and examples. While this analysis observes what developers include in documentation, it does not test which instructions actually improve generated code. This study addresses how to close that gap by testing specifications through generated implementations and using the failures to guide refinement.

ClarifyGPT provides evidence that refining requirements can improve code generation \citep{mu2024}. It compares generated solutions to identify ambiguity, asks targeted clarification questions, and incorporates the answers into revised requirements. In an evaluation with ten participants, GPT-4's rate of generating a solution that passes the tests on its first attempt (Pass@1) increases from 70.96\% to 80.80\% on a manually verified set of basic Python programming tasks (MBPP-sanitized). This supports refinement as a way to improve functional correctness, but does not prove its effectiveness for improving architectural implementation quality. This study instead uses reference code and evaluation feedback to identify what the specification fails to communicate.

Refinement also requires deciding how to organize a software specification. \citet{feng2026} compare natural language descriptions, Gherkin scenarios, domain models, and API signature models for MVC business logic generation across three systems and five LLMs. Adding structured models improves overall test performance, although Gherkin does not consistently outperform natural language when both include models. Their results support that supplying explicit, structured design information can improve LLM output quality. While their experiment uses fixed inputs without a feedback loop, this study revises specification content and organization in response to implementation failures.

Evaluating those failures requires assessing properties that tests alone do not establish. \citet{wang2025} compare LLM judging methods with human assessments of code generation, translation, and summarization. They find high correlations with human scores for generation and translation. This supports using LLM judgments for evaluating implementation quality not only on functional correctness but architectural adherence as well.

\section{Methods}

\subsection{Research Design}

In this experiment, we determine whether revising specifications improves LLM-generated code. For each FEA-Bench feature, ten agents generate implementations from the PR description. An editor revises the specification using evaluator feedback and reference code. Ten fresh agents test each revision. Using ten independent generations for each iteration gives greater confidence in the strength of the specification despite the non-determinism of LLMs. Across iterations, the original repository, model settings, and evaluation criteria remain unchanged. We compare how many implementations pass every criterion before and after refinement. This comparison tests whether the refined specification improves results over the original PR description and which changes are associated with those improvements. After a feature's specification converges, we simplify it. We repeatedly shorten the specification and keep the shortest version that still passes evaluation to examine which instructions the implementations actually need.

\subsection{Dataset and Preparation}

We use all 1,401 feature tasks from the 83 repositories in FEA-Bench. Each task includes a PR describing a feature and the merged code, which serves as the reference implementation. FEA-Bench was chosen since it focuses on adding features to existing repositories rather than solving standalone problems, which is the scope of this study. To prepare the task, a script checks out the repository as it was before the feature was added and installs the dependencies needed to run its tests. Each agent then starts with a fresh copy of this original code. If we cannot prepare the environment or a run is interrupted, we record that separately from an implementation that fails evaluation.

\subsection{Iterative Specification Refinement}

The baseline specification used to generate the first implementations is just the original PR title and description, including any code they contain. Within the FEA-Bench dataset, some PR descriptions differ from the merged implementation due to feature requirements changing in response to PR comments. We therefore use the merged code as the source of truth for feature intent. Each of the ten code-generating agents receives this baseline specification and a copy of the original repository in its own workspace. The agents can inspect the repository and run commands as they implement the feature, but they do not receive the reference code or feedback from other runs. All agents, including the evaluator and editor, use GPT-6 Luna with high reasoning effort through headless Codex. The reason headless Codex was chosen was so that the agent could inspect the repository and run commands without manual interaction. We chose Luna as a low-cost model because our goal is to make specifications clear enough for even an inexpensive model to implement them. We keep the model settings and tools the same throughout the loop and do not impose a token budget.

After an agent finishes, the script checks that its patch applies and runs the configured tests. These include tests the agent changed and regression tests from the reference implementation. However, passing tests alone does not establish that the code follows the repository's design. We chose the following eleven criteria to cover both functional correctness and integration with the existing repository. The evaluator therefore reviews the generated code and test results against eleven criteria, marking each as pass or fail with code evidence:

\begin{enumerate}
  \item Patch application and executable checks
  \item File placement
  \item API contracts
  \item Dependency direction
  \item Required registration and interface implementation
  \item Approved dependencies
  \item Architectural fit
  \item Use of existing extension mechanisms
  \item Separation of responsibilities across architectural boundaries
  \item Repository conventions
  \item Requirement intent
\end{enumerate}

An implementation passes only if it passes all eleven criteria.

To make these judgments, the evaluator has access to both the original repository and the reference implementation. It uses the reference to determine what the feature should do and how it should fit the repository. Importantly, a line-by-line match isn’t needed for the implementation to pass. Also, we do not give the evaluator the specification or PR description as requirements. Otherwise, it could accept code that follows an incorrect specification instead of implementing the intended feature.

When the evaluator agent identifies failures, the editor agent receives this feedback and the specification used to generate the failed code. It also has access to the generated code and both repository states so it can investigate why an implementation failed. Based on this information, it produces one revised specification. The editor can then add, remove, or modify sections, but any section it adds must come from our fixed list of specification categories. These categories were chosen as common specification categories that cover a feature's purpose, behavior, and implementation constraints. Keeping the available headings consistent across features lets us compare what headings the editor uses across features. The headings of the allowed sections are organized as follows.

\begin{table}[t]
\caption{Specification categories available to the editor.
Individual headings, rather than groups, form specification sections.}
\label{tab:specification-categories}
\small
\begin{tabular}{@{}p{0.23\linewidth}p{0.73\linewidth}@{}}
\toprule
Group & Allowed headings \\
\midrule
Purpose & Goals; Scope; Assumptions; Terminology \\
Behavior & Workflows \\
Data & Data model \\
Interfaces & API contracts \\
Architecture & Module architecture; Implementation constraints;
Dependencies \\
Conventions & Naming; Validation; Error handling; Logging;
Documentation \\
Quality & Performance; Security; Reliability; Usability;
Maintainability \\
Compatibility & Compatibility \\
Operations & Operations \\
Acceptance & Behavioral acceptance criteria; Test requirements \\
\bottomrule
\end{tabular}
\end{table}

Within these sections, the specification can name APIs and describe their contracts, but it cannot provide actual code because this would blur the line between software specifications and software implementation. We include the category list and editor prompt with the experiment records.

Ten fresh agents generate code from each revised specification and the original repository. We call the specification converged only if all ten implementations pass every criterion. We require every sampled implementation to pass because we want a specification that works consistently across attempts. If any agent fails, we start a new iteration using the failed evaluations to revise the specification. If a feature does not converge within ten revisions beyond the baseline, we stop the loop and record it as non-convergent. This bounds the work spent on a feature that doesn't converge.

[Insert a figure showing the iteration loop and inputs to the agents.]

\subsection{Collected Data and Analysis}

We save each specification with its generated code diffs and evaluation results. By comparing these records between revisions, we can see which instruction changes accompany which code changes. We compare the proportion of implementations passing each evaluation criterion before and after refinement, and report how many features converge and how many revisions they require.

To explain changes in criterion pass rates across revisions, we examine how the instructions change after specific failures. In the Seaborn experiment, all ten implementations from the first revision passed executable checks. However, four still failed convention checks involving documentation or test organization. In response, the editor clarified those requirements, and all ten implementations in the next batch passed. We look for similar cases across the dataset and compare the instructions before and after each correction. The limit to this is that if several instructions change together, we cannot say that one of them alone caused the improvement. However, analyzing the instructions across the entire FEA-bench dataset reveals common trends that we can observe.

\subsection{Specification Simplification}

A converged specification can still contain instructions that the implementation does not need, since the editor only adds or clarifies instructions in response to failures. After the loop, we therefore simplify each converged specification. In each round, three independent editor sessions shorten the current specification by 10--30\% of its words. Each session receives the same specification, evaluation feedback, and repositories as the editor, but cannot see the other sessions' drafts. We ask them to first remove repeated instructions and explanations that do not help implement the feature, then express the remaining requirements more directly. The same rules apply as during refinement: the sessions can only use the approved headings and cannot provide actual code.

Ten fresh agents generate code from each draft, using the same generation and evaluation settings as the loop. Among the drafts for which all ten implementations pass every criterion, we keep the shortest as the specification for the next round. If no draft passes, we stop and keep the last passing specification. We bound each reduction to 10--30\% so that a round removes enough content to be meaningful without discarding too much at once, and we generate three drafts because a single failed attempt does not mean that the specification cannot be shortened further.

We compare each simplified specification with the converged version to see which sections and instructions were removed and which remained. However, a passing draft does not establish that every removed instruction was unnecessary individually, since several instructions are usually removed together and some may also appear elsewhere in the specification. A draft can also fail because of variation in generation or an evaluator error. We therefore use these results to identify instructions worth investigating, rather than conclude that a category is always necessary or unnecessary.

\subsection{Guidance for New Repositories}

We use these analyses to identify possible guidance for writing specifications for new features. For example, if generated code repeatedly fails to register a new component, the guidance could ask a specification-writing agent to find the repository's registration mechanism and explain how to use it. The instructions that remain after simplification then help us decide which instructions to investigate further. The goal is to package this guidance as a reusable skill and test it on new repositories without access to their completed implementations. However, the procedure for developing and testing that skill remains outside the current experiment.

\subsection{Validity and Reproducibility}

We manually evaluate a random sample of patches using the same criteria as the evaluator. We compare our decisions with its judgments and report where it accepts an implementation we reject, or rejects one we accept. We do not manually correct individual judgments during the loop. The generator is instructed not to read reference code or results from other runs, but all agents have full filesystem access. These restrictions therefore depend on the agents following their prompts; network access is disabled. Model variation and evaluator errors can affect the results. The reference code may also contain flaws, and repeated instructions can make a section appear less useful when removed. We retain the prompts and run settings alongside the experiment artifacts so others can repeat the procedure.

The environment in which a repository's code runs can also affect these evaluations. Mismatched dependencies and compatibility requirements may prevent code from running even when the implementation is correct, so environment failures must be distinguished from implementation failures. Furthermore, poor implementation of the testing environment can lead to information leakage between the ground truth implementation and the code-generation agent, for example. Deficiencies in the evaluation agent prompt may also lead to misleading judgments and poor convergence performance. These limitations matter because the refinement loop uses evaluation feedback to decide what instructions to change. Human oversight will be used in the experiments to mitigate these issues.

\section{Discussion}

\subsection{Limitations}

The reference code may also contain technical debt or depart from broader repository conventions.

\bibliographystyle{ACM-Reference-Format}
\bibliography{references}
\end{document}
